\documentclass[aspectratio=169, 11pt]{beamer}
% ============================================================================
% Preambulo compartido para presentaciones Beamer
% Estadistica 2026-II - Pontificia Universidad Javeriana
% Incluir con: % ============================================================================
% Preambulo compartido para presentaciones Beamer
% Estadistica 2026-II - Pontificia Universidad Javeriana
% Incluir con: % ============================================================================
% Preambulo compartido para presentaciones Beamer
% Estadistica 2026-II - Pontificia Universidad Javeriana
% Incluir con: \input{../preambulo_beamer.tex} despues de \documentclass
% ============================================================================

\usepackage[T1]{fontenc}
\usepackage[utf8]{inputenc}
\usepackage[spanish]{babel}
\usepackage{amsmath, amssymb, amsfonts}
\usepackage{booktabs}
\usepackage{graphicx}
\usepackage{tikz}
\usetikzlibrary{shapes.geometric, positioning, arrows.meta, calc}
\usepackage{pgfplots}
\pgfplotsset{compat=1.18}
\usepackage{listings}
\usepackage{xcolor}
\usepackage{hyperref}
\usepackage{multicol}

% --- Tema Beamer (minimalista) ---
\usetheme{default}
\usecolortheme{default}
\usefonttheme{default}
\setbeamertemplate{navigation symbols}{}
\setbeamertemplate{caption}[numbered]
\setbeamertemplate{footline}{%
  \hspace{1em}{\tiny\url{https://github.com/polanco-jaime/Curso-Estadistica-2026-3}}\hfill\usebeamerfont{footline}\insertframenumber/\inserttotalframenumber\hspace{1em}\vspace{0.5em}%
}
\setbeamertemplate{itemize items}[circle]
\setbeamertemplate{enumerate items}[default]

% --- Colores institucionales (sutiles) ---
\definecolor{javeriana}{RGB}{0, 62, 105}
\definecolor{javerianalight}{RGB}{0, 105, 170}
\definecolor{codigobg}{RGB}{245, 245, 245}
\definecolor{codigofg}{RGB}{40, 40, 40}
\definecolor{comentario}{RGB}{100, 100, 100}
\definecolor{keyword}{RGB}{0, 100, 150}
\definecolor{string}{RGB}{180, 60, 30}

\setbeamercolor{title}{fg=javeriana}
\setbeamercolor{frametitle}{fg=javeriana}
\setbeamercolor{structure}{fg=javeriana}
\setbeamercolor{block title}{fg=white, bg=javeriana!75}
\setbeamercolor{block body}{bg=javeriana!5}
\setbeamercolor{block title alerted}{fg=white, bg=red!70!black}
\setbeamercolor{block body alerted}{bg=red!5}
\setbeamercolor{block title example}{fg=white, bg=green!40!black}
\setbeamercolor{block body example}{bg=green!5}
\setbeamercolor{item}{fg=javeriana}

\setbeamerfont{title}{series=\bfseries, size=\Large}
\setbeamerfont{frametitle}{series=\bfseries}

% --- Configuracion de listings para R ---
\lstset{
  language=R,
  basicstyle=\ttfamily\footnotesize,
  backgroundcolor=\color{codigobg},
  commentstyle=\color{comentario}\itshape,
  keywordstyle=\color{keyword}\bfseries,
  stringstyle=\color{string},
  numbers=none,
  frame=single,
  rulecolor=\color{javeriana!30},
  breaklines=true,
  showstringspaces=false,
  columns=flexible,
  morekeywords={library, ggplot, aes, geom_histogram, geom_boxplot,
    geom_point, geom_smooth, geom_density, geom_bar, geom_line,
    geom_vline, geom_hline, geom_errorbar, labs, theme_minimal,
    facet_wrap, scale_fill_manual, scale_color_manual,
    summarise, group_by, filter, mutate, select, arrange,
    read.csv, head, str, summary, mean, median, sd, var, cor,
    lm, t.test, chisq.test, wilcox.test, kruskal.test, aov,
    pnorm, qnorm, dnorm, rnorm, qt, pt, qchisq, pchisq, qf, pf,
    confint, coeftest, vcovHC, bptest, vif, cooks.distance,
    TukeyHSD, table, prop.table}
}

% --- Comandos personalizados ---
\newcommand{\R}{\texttt{R}}
\newcommand{\code}[1]{\texttt{\footnotesize #1}}
\newcommand{\variable}[1]{\texttt{#1}}
\newcommand{\paquete}[1]{\texttt{#1}}

% --- Informacion del curso ---
\institute{Pontificia Universidad Javeriana\\
  Facultad de Ciencias Econ\'omicas y Administrativas}
\date{2026-II}
 despues de \documentclass
% ============================================================================

\usepackage[T1]{fontenc}
\usepackage[utf8]{inputenc}
\usepackage[spanish]{babel}
\usepackage{amsmath, amssymb, amsfonts}
\usepackage{booktabs}
\usepackage{graphicx}
\usepackage{tikz}
\usetikzlibrary{shapes.geometric, positioning, arrows.meta, calc}
\usepackage{pgfplots}
\pgfplotsset{compat=1.18}
\usepackage{listings}
\usepackage{xcolor}
\usepackage{hyperref}
\usepackage{multicol}

% --- Tema Beamer (minimalista) ---
\usetheme{default}
\usecolortheme{default}
\usefonttheme{default}
\setbeamertemplate{navigation symbols}{}
\setbeamertemplate{caption}[numbered]
\setbeamertemplate{footline}{%
  \hspace{1em}{\tiny\url{https://github.com/polanco-jaime/Curso-Estadistica-2026-3}}\hfill\usebeamerfont{footline}\insertframenumber/\inserttotalframenumber\hspace{1em}\vspace{0.5em}%
}
\setbeamertemplate{itemize items}[circle]
\setbeamertemplate{enumerate items}[default]

% --- Colores institucionales (sutiles) ---
\definecolor{javeriana}{RGB}{0, 62, 105}
\definecolor{javerianalight}{RGB}{0, 105, 170}
\definecolor{codigobg}{RGB}{245, 245, 245}
\definecolor{codigofg}{RGB}{40, 40, 40}
\definecolor{comentario}{RGB}{100, 100, 100}
\definecolor{keyword}{RGB}{0, 100, 150}
\definecolor{string}{RGB}{180, 60, 30}

\setbeamercolor{title}{fg=javeriana}
\setbeamercolor{frametitle}{fg=javeriana}
\setbeamercolor{structure}{fg=javeriana}
\setbeamercolor{block title}{fg=white, bg=javeriana!75}
\setbeamercolor{block body}{bg=javeriana!5}
\setbeamercolor{block title alerted}{fg=white, bg=red!70!black}
\setbeamercolor{block body alerted}{bg=red!5}
\setbeamercolor{block title example}{fg=white, bg=green!40!black}
\setbeamercolor{block body example}{bg=green!5}
\setbeamercolor{item}{fg=javeriana}

\setbeamerfont{title}{series=\bfseries, size=\Large}
\setbeamerfont{frametitle}{series=\bfseries}

% --- Configuracion de listings para R ---
\lstset{
  language=R,
  basicstyle=\ttfamily\footnotesize,
  backgroundcolor=\color{codigobg},
  commentstyle=\color{comentario}\itshape,
  keywordstyle=\color{keyword}\bfseries,
  stringstyle=\color{string},
  numbers=none,
  frame=single,
  rulecolor=\color{javeriana!30},
  breaklines=true,
  showstringspaces=false,
  columns=flexible,
  morekeywords={library, ggplot, aes, geom_histogram, geom_boxplot,
    geom_point, geom_smooth, geom_density, geom_bar, geom_line,
    geom_vline, geom_hline, geom_errorbar, labs, theme_minimal,
    facet_wrap, scale_fill_manual, scale_color_manual,
    summarise, group_by, filter, mutate, select, arrange,
    read.csv, head, str, summary, mean, median, sd, var, cor,
    lm, t.test, chisq.test, wilcox.test, kruskal.test, aov,
    pnorm, qnorm, dnorm, rnorm, qt, pt, qchisq, pchisq, qf, pf,
    confint, coeftest, vcovHC, bptest, vif, cooks.distance,
    TukeyHSD, table, prop.table}
}

% --- Comandos personalizados ---
\newcommand{\R}{\texttt{R}}
\newcommand{\code}[1]{\texttt{\footnotesize #1}}
\newcommand{\variable}[1]{\texttt{#1}}
\newcommand{\paquete}[1]{\texttt{#1}}

% --- Informacion del curso ---
\institute{Pontificia Universidad Javeriana\\
  Facultad de Ciencias Econ\'omicas y Administrativas}
\date{2026-II}
 despues de \documentclass
% ============================================================================

\usepackage[T1]{fontenc}
\usepackage[utf8]{inputenc}
\usepackage[spanish]{babel}
\usepackage{amsmath, amssymb, amsfonts}
\usepackage{booktabs}
\usepackage{graphicx}
\usepackage{tikz}
\usetikzlibrary{shapes.geometric, positioning, arrows.meta, calc}
\usepackage{pgfplots}
\pgfplotsset{compat=1.18}
\usepackage{listings}
\usepackage{xcolor}
\usepackage{hyperref}
\usepackage{multicol}

% --- Tema Beamer (minimalista) ---
\usetheme{default}
\usecolortheme{default}
\usefonttheme{default}
\setbeamertemplate{navigation symbols}{}
\setbeamertemplate{caption}[numbered]
\setbeamertemplate{footline}{%
  \hspace{1em}{\tiny\url{https://github.com/polanco-jaime/Curso-Estadistica-2026-3}}\hfill\usebeamerfont{footline}\insertframenumber/\inserttotalframenumber\hspace{1em}\vspace{0.5em}%
}
\setbeamertemplate{itemize items}[circle]
\setbeamertemplate{enumerate items}[default]

% --- Colores institucionales (sutiles) ---
\definecolor{javeriana}{RGB}{0, 62, 105}
\definecolor{javerianalight}{RGB}{0, 105, 170}
\definecolor{codigobg}{RGB}{245, 245, 245}
\definecolor{codigofg}{RGB}{40, 40, 40}
\definecolor{comentario}{RGB}{100, 100, 100}
\definecolor{keyword}{RGB}{0, 100, 150}
\definecolor{string}{RGB}{180, 60, 30}

\setbeamercolor{title}{fg=javeriana}
\setbeamercolor{frametitle}{fg=javeriana}
\setbeamercolor{structure}{fg=javeriana}
\setbeamercolor{block title}{fg=white, bg=javeriana!75}
\setbeamercolor{block body}{bg=javeriana!5}
\setbeamercolor{block title alerted}{fg=white, bg=red!70!black}
\setbeamercolor{block body alerted}{bg=red!5}
\setbeamercolor{block title example}{fg=white, bg=green!40!black}
\setbeamercolor{block body example}{bg=green!5}
\setbeamercolor{item}{fg=javeriana}

\setbeamerfont{title}{series=\bfseries, size=\Large}
\setbeamerfont{frametitle}{series=\bfseries}

% --- Configuracion de listings para R ---
\lstset{
  language=R,
  basicstyle=\ttfamily\footnotesize,
  backgroundcolor=\color{codigobg},
  commentstyle=\color{comentario}\itshape,
  keywordstyle=\color{keyword}\bfseries,
  stringstyle=\color{string},
  numbers=none,
  frame=single,
  rulecolor=\color{javeriana!30},
  breaklines=true,
  showstringspaces=false,
  columns=flexible,
  morekeywords={library, ggplot, aes, geom_histogram, geom_boxplot,
    geom_point, geom_smooth, geom_density, geom_bar, geom_line,
    geom_vline, geom_hline, geom_errorbar, labs, theme_minimal,
    facet_wrap, scale_fill_manual, scale_color_manual,
    summarise, group_by, filter, mutate, select, arrange,
    read.csv, head, str, summary, mean, median, sd, var, cor,
    lm, t.test, chisq.test, wilcox.test, kruskal.test, aov,
    pnorm, qnorm, dnorm, rnorm, qt, pt, qchisq, pchisq, qf, pf,
    confint, coeftest, vcovHC, bptest, vif, cooks.distance,
    TukeyHSD, table, prop.table}
}

% --- Comandos personalizados ---
\newcommand{\R}{\texttt{R}}
\newcommand{\code}[1]{\texttt{\footnotesize #1}}
\newcommand{\variable}[1]{\texttt{#1}}
\newcommand{\paquete}[1]{\texttt{#1}}

% --- Informacion del curso ---
\institute{Pontificia Universidad Javeriana\\
  Facultad de Ciencias Econ\'omicas y Administrativas}
\date{2026-II}

\title{Sesión 12: Integración y Análisis Aplicado}
\subtitle{El pipeline estadístico completo con datos ICFES}
\author{Prof.\ Jaime Polanco Jim\'enez}
\date{Viernes 13 de noviembre de 2026}

\begin{document}

\begin{frame}
\titlepage
\end{frame}

\begin{frame}{Agenda}
\tableofcontents
\end{frame}

% ============================================================================
\section{El pipeline estadístico}
% ============================================================================

\begin{frame}{El pipeline estadístico completo}
\begin{center}
\begin{tikzpicture}[
  node distance=0.8cm and 1.2cm,
  box/.style={rectangle, draw=javeriana, fill=javeriana!10, thick, minimum width=2.2cm, minimum height=0.7cm, rounded corners, font=\small},
  arrow/.style={->, thick, javeriana}
]
\node[box] (pregunta) {Pregunta};
\node[box, below=of pregunta] (datos) {Datos};
\node[box, below=of datos] (eda) {EDA};
\node[box, below=of eda] (ic) {IC};
\node[box, right=of pregunta] (pruebas) {Pruebas};
\node[box, below=of pruebas] (modelo) {Modelo};
\node[box, below=of modelo] (diag) {Diagnósticos};
\node[box, below=of diag] (concl) {Conclusiones};

\draw[arrow] (pregunta) -- (datos);
\draw[arrow] (datos) -- (eda);
\draw[arrow] (eda) -- (ic);
\draw[arrow] (ic) -- (pruebas);
\draw[arrow] (pruebas) -- (modelo);
\draw[arrow] (modelo) -- (diag);
\draw[arrow] (diag) -- (concl);

% Feedback loops
\draw[arrow, dashed, bend left=20] (diag.west) to (modelo.west);
\draw[arrow, dashed, bend left=30] (concl.east) to (pregunta.east);
\end{tikzpicture}
\end{center}

\vspace{0.3cm}
\begin{block}{El análisis estadístico es un proceso iterativo}
\begin{itemize}
\item Cada paso informa el siguiente
\item Los diagnósticos pueden llevarnos a re-especificar el modelo
\item Las conclusiones pueden generar nuevas preguntas
\item Los datos cuentan una historia — nuestro trabajo es descubrirla con rigor
\end{itemize}
\end{block}
\end{frame}

\begin{frame}{¿Por qué importa el pipeline completo?}
\begin{columns}
\begin{column}{0.48\textwidth}
\begin{alertblock}{Mal análisis}
\begin{itemize}
\item Saltar directamente a regresión
\item No explorar los datos
\item Ignorar supuestos
\item No verificar diagnósticos
\item Reportar solo p-valores
\end{itemize}

\vspace{0.3cm}
\textbf{Resultado:}
\begin{itemize}
\item Conclusiones incorrectas
\item Modelos mal especificados
\item Inferencia inválida
\end{itemize}
\end{alertblock}
\end{column}

\begin{column}{0.48\textwidth}
\begin{exampleblock}{Buen análisis}
\begin{itemize}
\item Comenzar con una pregunta clara
\item Explorar antes de modelar
\item Verificar supuestos
\item Diagnosticar problemas
\item Reportar efecto + incertidumbre
\end{itemize}

\vspace{0.3cm}
\textbf{Resultado:}
\begin{itemize}
\item Conclusiones sólidas
\item Modelos justificados
\item Inferencia válida
\end{itemize}
\end{exampleblock}
\end{column}
\end{columns}

\vspace{0.3cm}
\begin{center}
\textbf{Objetivo de hoy}: aplicar el pipeline completo a datos ICFES de Negocios Internacionales.
\end{center}
\end{frame}

\begin{frame}{Caso de estudio: competencia en inglés en NI}
\begin{block}{Contexto}
Los programas de Negocios Internacionales requieren dominio del inglés como competencia clave. El examen Saber Pro evalúa esta competencia con la prueba de inglés (MOD\_INGLES\_PUNT, escala 0-300, convertible a niveles MCER: A-, A1, A2, B1, B2).
\end{block}

\vspace{0.3cm}
\textbf{Pregunta de investigación:}
\begin{center}
\large
\textcolor{javeriana}{¿Qué factores socioeconómicos y académicos determinan\\el desempeño en inglés de estudiantes de Negocios Internacionales\\en Colombia?}
\end{center}

\vspace{0.3cm}
\textbf{Relevancia:}
\begin{itemize}
\item Informar políticas de admisión y nivelación
\item Identificar grupos vulnerables que requieren apoyo
\item Evaluar la equidad en el acceso a formación en inglés
\end{itemize}
\end{frame}

% ============================================================================
\section{Paso 1: Pregunta de investigación}
% ============================================================================

\begin{frame}{Formulando una buena pregunta de investigación}
\begin{block}{Características de una buena pregunta}
\begin{itemize}
\item \textbf{Clara}: términos bien definidos
\item \textbf{Medible}: variables observables
\item \textbf{Relevante}: tiene implicaciones prácticas o teóricas
\item \textbf{Acotada}: no demasiado amplia ni demasiado estrecha
\end{itemize}
\end{block}

\vspace{0.3cm}
\begin{columns}
\begin{column}{0.48\textwidth}
\begin{alertblock}{Pregunta vaga}
"¿Por qué algunos estudiantes son mejores en inglés?"

\vspace{0.2cm}
\small
\textcolor{red}{Problemas:}
\begin{itemize}
\item ¿Qué es "mejor"?
\item ¿Qué estudiantes?
\item ¿Qué factores?
\end{itemize}
\end{alertblock}
\end{column}

\begin{column}{0.48\textwidth}
\begin{exampleblock}{Pregunta específica}
"¿Cuál es el efecto del estrato socioeconómico y del tipo de IES en el puntaje de inglés en Saber Pro para estudiantes de NI?"

\vspace{0.2cm}
\small
\textcolor{javerianalight}{Bien definida:}
\begin{itemize}
\item Variable dependiente: puntaje inglés
\item Variables independientes: estrato, tipo IES
\item Población: estudiantes NI
\end{itemize}
\end{exampleblock}
\end{column}
\end{columns}
\end{frame}

\begin{frame}{Operacionalización de variables}
\begin{center}
\renewcommand{\arraystretch}{1.5}
\begin{tabular}{lll}
\toprule
\textbf{Concepto} & \textbf{Variable} & \textbf{Medición} \\
\midrule
Competencia en inglés & \variable{MOD\_INGLES\_PUNT} & Escala 0-300 (continua) \\
 & \variable{NIVEL\_MCER} & A-, A1, A2, B1, B2 (ordinal) \\
\midrule
Nivel socioeconómico & \variable{ESTU\_ESTRATO} & 1-6 (ordinal) \\
\midrule
Tipo de institución & \variable{PRIVADA} & 0 = pública, 1 = privada \\
\midrule
Desempeño previo & \variable{PUNT\_INGLES\_S11} & Puntaje Saber 11 (continua) \\
\midrule
Género & \variable{ESTU\_GENERO} & M / F (binaria) \\
\midrule
Región & \variable{REGION} & Andina, Caribe, etc. (categórica) \\
\bottomrule
\end{tabular}
\end{center}

\vspace{0.3cm}
\small
\textbf{Nota}: La operacionalización conecta conceptos abstractos con mediciones concretas.
\end{frame}

% ============================================================================
\section{Paso 2: Cargar y limpiar datos}
% ============================================================================

\begin{frame}[fragile]{Cargando los datos}
\begin{lstlisting}
# Cargar datos ICFES
icfes <- read.csv("icfes_saber_pro_2023.csv")

# Inspeccion inicial
str(icfes)      # Estructura de los datos
head(icfes)     # Primeras 6 filas
dim(icfes)      # Dimensiones (filas x columnas)
names(icfes)    # Nombres de las columnas

# Ver estadisticos basicos
summary(icfes)
\end{lstlisting}

\vspace{0.3cm}
\begin{exampleblock}{Output esperado}
\small
\begin{verbatim}
'data.frame':  234567 obs. of  42 variables:
 $ ESTU_CONSECUTIVO    : int  1 2 3 4 5 ...
 $ PROGRAMA            : chr "MEDICINA" "DERECHO" "INGENIERIA CIVIL" ...
 $ MOD_INGLES_PUNT     : num  165 142 178 NA 155 ...
 $ ESTU_ESTRATO        : int  3 2 4 1 3 ...
 ...
\end{verbatim}
\end{exampleblock}
\end{frame}

\begin{frame}[fragile]{Filtrando y creando variables}
\begin{lstlisting}
library(dplyr)

# Filtrar solo estudiantes de Negocios Internacionales
ni_data <- icfes %>%
  filter(PROGRAMA == "NEGOCIOS INTERNACIONALES")

# Ver cuantos estudiantes quedaron
nrow(ni_data)

# Crear variable PRIVADA
ni_data <- ni_data %>%
  mutate(PRIVADA = ifelse(TIPO_IES == "PRIVADA", 1, 0))

# Crear variable REGION
ni_data <- ni_data %>%
  mutate(REGION = case_when(
    DEPTO %in% c("CUNDINAMARCA", "BOYACA", "SANTANDER") ~ "Andina",
    DEPTO %in% c("ATLANTICO", "BOLIVAR", "MAGDALENA") ~ "Caribe",
    DEPTO %in% c("VALLE", "CAUCA", "NARINO") ~ "Pacifico",
    DEPTO %in% c("META", "CASANARE") ~ "Orinoquia",
    TRUE ~ "Otra"
  ))
\end{lstlisting}
\end{frame}

\begin{frame}[fragile]{Tratando valores faltantes}
\begin{lstlisting}
# Identificar valores faltantes
summary(ni_data$MOD_INGLES_PUNT)
sum(is.na(ni_data$MOD_INGLES_PUNT))

# Ver patron de valores faltantes
colSums(is.na(ni_data))

# Eliminar filas con NA en variables clave
ni_data <- ni_data %>%
  filter(!is.na(MOD_INGLES_PUNT),
         !is.na(ESTU_ESTRATO),
         !is.na(PUNT_INGLES_S11))

# Ver cuantas observaciones quedaron
nrow(ni_data)
\end{lstlisting}

\vspace{0.3cm}
\begin{alertblock}{Importante}
\begin{itemize}
\item \code{na.omit()} elimina TODAS las filas con algún NA (puede ser demasiado drástico)
\item Mejor: eliminar NA solo en variables clave
\item Reportar cuántas observaciones se perdieron y por qué
\end{itemize}
\end{alertblock}
\end{frame}

\begin{frame}[fragile]{Verificando la limpieza}
\begin{lstlisting}
# Estadisticos descriptivos de variables clave
summary(ni_data %>% select(MOD_INGLES_PUNT, ESTU_ESTRATO,
                           PUNT_INGLES_S11, PRIVADA))

# Verificar rangos validos
# El puntaje de ingles debe estar entre 0 y 300
ni_data <- ni_data %>%
  filter(MOD_INGLES_PUNT >= 0 & MOD_INGLES_PUNT <= 300)

# Verificar que el estrato este entre 1 y 6
ni_data <- ni_data %>%
  filter(ESTU_ESTRATO >= 1 & ESTU_ESTRATO <= 6)

# Crear tabla de frecuencias para variables categoricas
table(ni_data$PRIVADA)
table(ni_data$REGION)
table(ni_data$ESTU_GENERO)
\end{lstlisting}

\vspace{0.3cm}
\small
\textbf{Siempre verificar:} rangos válidos, categorías correctas, no haya valores imposibles.
\end{frame}

% ============================================================================
\section{Paso 3: EDA}
% ============================================================================

\begin{frame}[fragile]{EDA: estadísticos descriptivos}
\begin{lstlisting}
library(dplyr)

# Estadisticos basicos para MOD_INGLES_PUNT
ni_data %>%
  summarise(
    n = n(),
    media = mean(MOD_INGLES_PUNT),
    mediana = median(MOD_INGLES_PUNT),
    de = sd(MOD_INGLES_PUNT),
    minimo = min(MOD_INGLES_PUNT),
    maximo = max(MOD_INGLES_PUNT),
    q25 = quantile(MOD_INGLES_PUNT, 0.25),
    q75 = quantile(MOD_INGLES_PUNT, 0.75)
  )
\end{lstlisting}

\vspace{0.3cm}
\begin{exampleblock}{Output esperado}
\small
\begin{verbatim}
      n  media mediana    de minimo maximo   q25   q75
   2834 154.3   152.0  18.7   95.0  245.0 141.0 167.0
\end{verbatim}
\end{exampleblock}

\textbf{Interpretación:}
\begin{itemize}
\item Media $\approx$ mediana $\Rightarrow$ distribución aproximadamente simétrica
\item DE = 18.7 puntos (variabilidad moderada)
\item Rango: 95-245 (amplitud de 150 puntos)
\end{itemize}
\end{frame}

\begin{frame}[fragile]{EDA: desagregación por grupos}
\begin{lstlisting}
# Por estrato socioeconomico
ni_data %>%
  group_by(ESTU_ESTRATO) %>%
  summarise(
    n = n(),
    media = mean(MOD_INGLES_PUNT),
    de = sd(MOD_INGLES_PUNT)
  )

# Por tipo de IES
ni_data %>%
  group_by(PRIVADA) %>%
  summarise(
    n = n(),
    media = mean(MOD_INGLES_PUNT),
    de = sd(MOD_INGLES_PUNT)
  )

# Por region
ni_data %>%
  group_by(REGION) %>%
  summarise(
    n = n(),
    media = mean(MOD_INGLES_PUNT),
    de = sd(MOD_INGLES_PUNT)
  ) %>%
  arrange(desc(media))  # ordenar por media descendente
\end{lstlisting}
\end{frame}

\begin{frame}{EDA: hallazgos preliminares}
\begin{exampleblock}{Hallazgos esperados (ilustrativos)}
\small
\renewcommand{\arraystretch}{1.3}
\begin{tabular}{lrrr}
\toprule
\textbf{Estrato} & \textbf{n} & \textbf{Media} & \textbf{DE} \\
\midrule
1 & 234 & 138.5 & 20.3 \\
2 & 567 & 145.2 & 18.9 \\
3 & 891 & 152.8 & 17.5 \\
4 & 678 & 161.4 & 16.2 \\
5 & 345 & 169.7 & 15.8 \\
6 & 119 & 178.3 & 14.5 \\
\bottomrule
\end{tabular}
\end{exampleblock}

\textbf{Observaciones:}
\begin{itemize}
\item Relación positiva entre estrato y puntaje (gradiente claro)
\item La variabilidad (DE) disminuye al aumentar el estrato
\item Estudiantes de estrato 6 puntúan 40 puntos más que estrato 1 (diferencia sustantiva)
\end{itemize}

\vspace{0.3cm}
\small
\textcolor{javeriana}{Pregunta}: ¿Esta diferencia es estadísticamente significativa? $\to$ necesitamos pruebas formales.
\end{frame}

\begin{frame}[fragile]{EDA: visualizaciones}
\begin{lstlisting}
library(ggplot2)

# Histograma de puntaje de ingles
ggplot(ni_data, aes(x = MOD_INGLES_PUNT)) +
  geom_histogram(binwidth = 10, fill = "steelblue", color = "white") +
  labs(title = "Distribucion del puntaje de ingles en NI",
       x = "Puntaje ingles", y = "Frecuencia") +
  theme_minimal()

# Boxplot por estrato
ggplot(ni_data, aes(x = factor(ESTU_ESTRATO), y = MOD_INGLES_PUNT,
                    fill = factor(ESTU_ESTRATO))) +
  geom_boxplot() +
  labs(title = "Puntaje de ingles por estrato socioeconomico",
       x = "Estrato", y = "Puntaje ingles") +
  theme_minimal() +
  theme(legend.position = "none")
\end{lstlisting}
\end{frame}

\begin{frame}[fragile]{EDA: más visualizaciones}
\begin{lstlisting}
# Scatterplot: puntaje Saber 11 vs Saber Pro
ggplot(ni_data, aes(x = PUNT_INGLES_S11, y = MOD_INGLES_PUNT)) +
  geom_point(alpha = 0.3, color = "steelblue") +
  geom_smooth(method = "lm", se = TRUE, color = "red") +
  labs(title = "Relacion entre Saber 11 y Saber Pro (ingles)",
       x = "Puntaje ingles Saber 11",
       y = "Puntaje ingles Saber Pro") +
  theme_minimal()

# Barplot: proporcion que alcanza B1 por tipo de IES
ni_data <- ni_data %>%
  mutate(ALCANZA_B1 = MOD_INGLES_PUNT >= 169)

ni_data %>%
  group_by(PRIVADA) %>%
  summarise(prop_B1 = mean(ALCANZA_B1)) %>%
  ggplot(aes(x = factor(PRIVADA, labels = c("Publica", "Privada")),
             y = prop_B1, fill = factor(PRIVADA))) +
  geom_col() +
  scale_y_continuous(labels = scales::percent) +
  labs(title = "Proporcion que alcanza B1 por tipo de IES",
       x = "", y = "Proporcion") +
  theme_minimal() +
  theme(legend.position = "none")
\end{lstlisting}
\end{frame}

% ============================================================================
\section{Paso 4: Intervalos de confianza}
% ============================================================================

\begin{frame}[fragile]{IC para la media por región}
\begin{lstlisting}
library(dplyr)

# Calcular IC para cada region
ic_regiones <- ni_data %>%
  group_by(REGION) %>%
  summarise(
    n = n(),
    media = mean(MOD_INGLES_PUNT),
    de = sd(MOD_INGLES_PUNT),
    ee = de / sqrt(n),
    t_crit = qt(0.975, n - 1),
    lim_inf = media - t_crit * ee,
    lim_sup = media + t_crit * ee
  )

print(ic_regiones)
\end{lstlisting}

\vspace{0.3cm}
\begin{exampleblock}{Output esperado}
\tiny
\begin{verbatim}
  REGION        n  media    de    ee t_crit lim_inf lim_sup
  Andina     1234  156.2  17.8  0.51   1.96   155.2   157.2
  Caribe      567  148.5  19.3  0.81   1.96   146.9   150.1
  Pacifico    789  152.1  18.5  0.66   1.96   150.8   153.4
  Orinoquia   134  145.3  21.2  1.83   1.98   141.7   148.9
\end{verbatim}
\end{exampleblock}
\end{frame}

\begin{frame}[fragile]{Forest plot de IC por región}
\begin{lstlisting}
library(ggplot2)

ggplot(ic_regiones, aes(x = REGION, y = media)) +
  geom_point(size = 3, color = "steelblue") +
  geom_errorbar(aes(ymin = lim_inf, ymax = lim_sup), width = 0.2) +
  coord_flip() +
  labs(title = "Puntaje promedio de ingles por region (IC 95%)",
       x = "", y = "Puntaje promedio ingles") +
  theme_minimal()
\end{lstlisting}

\vspace{0.3cm}
\textbf{Interpretación del forest plot:}
\begin{itemize}
\item Barras horizontales = IC 95\%
\item Si dos IC NO se superponen, hay diferencia significativa
\item Región Andina tiene el puntaje más alto
\item Orinoquia tiene mayor incertidumbre (IC más amplio) debido a menor $n$
\end{itemize}
\end{frame}

\begin{frame}[fragile]{IC para proporciones}
\begin{lstlisting}
# Proporcion que alcanza B1 (>=169) por tipo de IES
ni_data %>%
  group_by(PRIVADA) %>%
  summarise(
    n = n(),
    exitos = sum(ALCANZA_B1),
    prop = mean(ALCANZA_B1),
    ee = sqrt(prop * (1 - prop) / n),
    z_crit = qnorm(0.975),
    lim_inf = prop - z_crit * ee,
    lim_sup = prop + z_crit * ee
  )
\end{lstlisting}

\vspace{0.3cm}
\begin{exampleblock}{Output esperado}
\small
\begin{verbatim}
  PRIVADA    n exitos  prop     ee z_crit lim_inf lim_sup
        0  987    234  0.237  0.014   1.96   0.210   0.264
        1 1847    658  0.356  0.011   1.96   0.334   0.378
\end{verbatim}
\end{exampleblock}

\textbf{Interpretación:}
Los IC no se superponen $\Rightarrow$ diferencia significativa: IES privadas tienen mayor proporción que alcanza B1 (35.6\% vs 23.7\%).
\end{frame}

% ============================================================================
\section{Paso 5: Pruebas de hipótesis}
% ============================================================================

\begin{frame}[fragile]{Prueba t: pública vs privada}
\begin{lstlisting}
# Test t de dos muestras
t.test(MOD_INGLES_PUNT ~ PRIVADA, data = ni_data, var.equal = FALSE)
\end{lstlisting}

\vspace{0.3cm}
\begin{exampleblock}{Output esperado}
\small
\begin{verbatim}
    Welch Two Sample t-test

data:  MOD_INGLES_PUNT by PRIVADA
t = -11.234, df = 1678.5, p-value < 2.2e-16
alternative hypothesis: true difference in means is not equal to 0
95 percent confidence interval:
 -13.56  -9.42
sample estimates:
mean in group 0  mean in group 1
          148.2            159.7
\end{verbatim}
\end{exampleblock}

\textbf{Conclusión:}
\begin{itemize}
\item $t = -11.23$, $p < 0.001$ $\Rightarrow$ diferencia altamente significativa
\item Estudiantes de IES privadas puntúan 11.5 puntos más en promedio
\item IC 95\%: [9.42, 13.56] puntos
\end{itemize}
\end{frame}

\begin{frame}[fragile]{ANOVA: diferencias entre regiones}
\begin{lstlisting}
# ANOVA de un factor
anova_regiones <- aov(MOD_INGLES_PUNT ~ REGION, data = ni_data)
summary(anova_regiones)

# Test post-hoc de Tukey
TukeyHSD(anova_regiones)
\end{lstlisting}

\vspace{0.3cm}
\begin{exampleblock}{Output ANOVA}
\small
\begin{verbatim}
              Df  Sum Sq Mean Sq F value   Pr(>F)
REGION         3   45678   15226   52.34 < 2e-16 ***
Residuals   2830  823456     291
\end{verbatim}
\end{exampleblock}

\textbf{Conclusión:}
\begin{itemize}
\item $F(3, 2830) = 52.34$, $p < 0.001$ $\Rightarrow$ hay diferencias significativas entre regiones
\item Tukey identifica qué pares difieren (ej: Andina vs Caribe: $p < 0.001$)
\end{itemize}
\end{frame}

\begin{frame}[fragile]{Chi-cuadrado: asociación nivel MCER $\times$ tipo IES}
\begin{lstlisting}
# Crear variable nivel MCER
ni_data <- ni_data %>%
  mutate(NIVEL_MCER = case_when(
    MOD_INGLES_PUNT < 142 ~ "A-",
    MOD_INGLES_PUNT < 169 ~ "A1-A2",
    TRUE ~ "B1+"
  ))

# Tabla de contingencia
tabla <- table(ni_data$NIVEL_MCER, ni_data$PRIVADA)
print(tabla)

# Test chi-cuadrado
chisq.test(tabla)
\end{lstlisting}

\vspace{0.3cm}
\begin{exampleblock}{Output chi-cuadrado}
\small
\begin{verbatim}
    Pearson's Chi-squared test
X-squared = 87.345, df = 2, p-value < 2.2e-16
\end{verbatim}
\end{exampleblock}

\textbf{Conclusión:} Hay asociación significativa entre nivel MCER y tipo de IES.
\end{frame}

\begin{frame}[fragile]{Tamaño del efecto: d de Cohen}
\begin{lstlisting}
library(effsize)

# d de Cohen para diferencia publica vs privada
cohen.d(MOD_INGLES_PUNT ~ PRIVADA, data = ni_data)
\end{lstlisting}

\vspace{0.3cm}
\begin{exampleblock}{Output}
\small
\begin{verbatim}
Cohen's d

d estimate: 0.62 (medium)
95 percent confidence interval:
    lower     upper
    0.51      0.73
\end{verbatim}
\end{exampleblock}

\textbf{Interpretación:}
\begin{itemize}
\item $d = 0.62$ (efecto mediano según Cohen: 0.2 = pequeño, 0.5 = mediano, 0.8 = grande)
\item La diferencia pública-privada es estadísticamente significativa Y sustantiva
\item Más allá del p-valor: el tamaño del efecto cuantifica la relevancia práctica
\end{itemize}
\end{frame}

% ============================================================================
\section{Paso 6: Modelo de regresión}
% ============================================================================

\begin{frame}[fragile]{Modelo simple: solo estrato}
\begin{lstlisting}
# Modelo simple
modelo1 <- lm(MOD_INGLES_PUNT ~ ESTU_ESTRATO, data = ni_data)
summary(modelo1)
\end{lstlisting}

\vspace{0.3cm}
\begin{exampleblock}{Output (ilustrativo)}
\tiny
\begin{verbatim}
Coefficients:
                 Estimate Std. Error t value Pr(>|t|)
(Intercept)       125.345      1.234   101.5   <2e-16 ***
ESTU_ESTRATO        7.234      0.345    20.9   <2e-16 ***

Residual standard error: 16.8 on 2832 degrees of freedom
Multiple R-squared:  0.134,	Adjusted R-squared:  0.133
\end{verbatim}
\end{exampleblock}

\textbf{Interpretación:}
\begin{itemize}
\item Cada punto adicional de estrato se asocia con 7.2 puntos más en inglés
\item $R^2 = 0.134$ $\Rightarrow$ el estrato explica 13.4\% de la varianza en inglés
\item Queda mucha varianza sin explicar (86.6\%)
\end{itemize}
\end{frame}

\begin{frame}[fragile]{Modelo múltiple: agregando variables}
\begin{lstlisting}
# Modelo multiple
modelo2 <- lm(MOD_INGLES_PUNT ~ ESTU_ESTRATO + PRIVADA +
              ESTU_GENERO + PUNT_INGLES_S11, data = ni_data)
summary(modelo2)
\end{lstlisting}

\vspace{0.3cm}
\begin{exampleblock}{Output (ilustrativo)}
\tiny
\begin{verbatim}
Coefficients:
                    Estimate Std. Error t value Pr(>|t|)
(Intercept)          45.234      3.456   13.08   <2e-16 ***
ESTU_ESTRATO          2.134      0.421    5.07  4.3e-07 ***
PRIVADA               5.432      1.234    4.40  1.1e-05 ***
ESTU_GENEROM         -1.543      0.987   -1.56    0.118
PUNT_INGLES_S11       0.876      0.034   25.76   <2e-16 ***

Residual standard error: 12.3 on 2829 degrees of freedom
Multiple R-squared:  0.428,	Adjusted R-squared:  0.427
\end{verbatim}
\end{exampleblock}

\textbf{Conclusiones:}
\begin{itemize}
\item Controlando por otras variables, el efecto del estrato se reduce (7.2 $\to$ 2.1)
\item El puntaje en Saber 11 es el predictor más fuerte ($t = 25.76$)
\item $R^2 = 0.428$ $\Rightarrow$ el modelo explica 42.8\% de la varianza
\end{itemize}
\end{frame}

\begin{frame}{Sesgo de variable omitida (OVB)}
\begin{block}{Comparación modelo1 vs modelo2}
\renewcommand{\arraystretch}{1.3}
\begin{tabular}{lrr}
\toprule
 & \textbf{Modelo 1} & \textbf{Modelo 2} \\
 & (solo estrato) & (múltiple) \\
\midrule
ESTU\_ESTRATO & 7.234*** & 2.134*** \\
PRIVADA & — & 5.432*** \\
PUNT\_INGLES\_S11 & — & 0.876*** \\
\midrule
$R^2$ & 0.134 & 0.428 \\
\bottomrule
\end{tabular}
\end{block}

\vspace{0.3cm}
\textbf{Explicación del cambio:}
\begin{itemize}
\item En modelo1, ESTU\_ESTRATO captura TANTO su efecto directo COMO el efecto indirecto de variables omitidas correlacionadas (PRIVADA, PUNT\_S11)
\item Estudiantes de estrato alto tienden a ir a IES privadas y tener mejor Saber 11
\item Al incluir estas variables en modelo2, el efecto de estrato se "purifica"
\item El efecto causal verdadero del estrato es ~2.1 puntos (controlado)
\end{itemize}
\end{frame}

\begin{frame}[fragile]{Interpretación de coeficientes}
\begin{lstlisting}
# Extraer coeficientes con IC
confint(modelo2)
\end{lstlisting}

\vspace{0.3cm}
\textbf{Interpretación de cada coeficiente (ceteris paribus):}
\begin{itemize}
\item \textbf{ESTU\_ESTRATO = 2.134}: subir un estrato se asocia con 2.1 puntos más en inglés, manteniendo constantes IES, género y Saber 11.
\item \textbf{PRIVADA = 5.432}: estudiar en IES privada se asocia con 5.4 puntos más, manteniendo constantes estrato, género y Saber 11.
\item \textbf{PUNT\_INGLES\_S11 = 0.876}: cada punto adicional en Saber 11 se asocia con 0.88 puntos más en Saber Pro.
\item \textbf{ESTU\_GENERO = -1.543} (no significativo, $p = 0.118$): hombres puntúan 1.5 puntos menos que mujeres, pero la diferencia no es estadísticamente significativa.
\end{itemize}

\vspace{0.3cm}
\small
\textcolor{javeriana}{Recordar}: estos son efectos \textit{parciales}, no causales (datos observacionales).
\end{frame}

% ============================================================================
\section{Paso 7: Diagnósticos}
% ============================================================================

\begin{frame}[fragile]{Diagnósticos gráficos}
\begin{lstlisting}
# Los 4 graficos diagnosticos
par(mfrow = c(2, 2))
plot(modelo2)
par(mfrow = c(1, 1))
\end{lstlisting}

\textbf{Revisión visual:}
\begin{enumerate}
\item \textbf{Residuals vs Fitted}: ¿Hay patrón? ¿Embudo?
  \begin{itemize}
  \item Buscar: línea roja horizontal, puntos dispersos aleatoriamente
  \end{itemize}
\item \textbf{Q-Q plot}: ¿Normalidad de residuos?
  \begin{itemize}
  \item Buscar: puntos cerca de la línea diagonal
  \end{itemize}
\item \textbf{Scale-Location}: ¿Homocedasticidad?
  \begin{itemize}
  \item Buscar: línea roja horizontal
  \end{itemize}
\item \textbf{Residuals vs Leverage}: ¿Observaciones influyentes?
  \begin{itemize}
  \item Buscar: puntos dentro de las líneas de Cook
  \end{itemize}
\end{enumerate}
\end{frame}

\begin{frame}[fragile]{Test de Breusch-Pagan}
\begin{lstlisting}
library(lmtest)

# Test de heterocedasticidad
bptest(modelo2)
\end{lstlisting}

\vspace{0.3cm}
\begin{exampleblock}{Output esperado}
\small
\begin{verbatim}
    studentized Breusch-Pagan test

data:  modelo2
BP = 38.472, df = 4, p-value = 8.741e-08
\end{verbatim}
\end{exampleblock}

\textbf{Conclusión:}
\begin{itemize}
\item $p < 0.001$ $\Rightarrow$ rechazamos $H_0$ de homocedasticidad
\item Hay evidencia de heterocedasticidad
\item Solución: usar errores estándar robustos
\end{itemize}
\end{frame}

\begin{frame}[fragile]{Observaciones influyentes}
\begin{lstlisting}
# Distancia de Cook
cook_d <- cooks.distance(modelo2)

# Umbral
umbral <- 4 / nrow(ni_data)

# Identificar influyentes
influyentes <- which(cook_d > umbral)
length(influyentes)

# Ver las 5 observaciones mas influyentes
ni_data[head(order(cook_d, decreasing = TRUE), 5), ]

# Graficar
plot(cook_d, type = "h", main = "Distancia de Cook")
abline(h = umbral, col = "red", lty = 2)
\end{lstlisting}

\vspace{0.3cm}
\small
\textbf{Acción}: si hay observaciones con $D_i > 4/n$, investigar (¿error de datos? ¿caso especial?). Si son legítimas, reportar análisis de sensibilidad.
\end{frame}

% ============================================================================
\section{Paso 8: Conclusiones y recomendaciones}
% ============================================================================

\begin{frame}{Síntesis de hallazgos}
\begin{block}{Principales resultados del análisis}
\begin{enumerate}
\item \textbf{Desempeño general}: El puntaje promedio de inglés en NI es 154.3 (DE = 18.7), equivalente a nivel A2/B1 en el MCER.

\item \textbf{Gradiente socioeconómico}: Existe una relación positiva entre estrato y puntaje. Controlando por otras variables, cada estrato adicional se asocia con 2.1 puntos más (IC 95\%: [1.3, 2.9]).

\item \textbf{Brecha público-privada}: Estudiantes de IES privadas puntúan 5.4 puntos más que los de públicas (IC 95\%: [3.0, 7.8]), controlando por estrato y desempeño previo.

\item \textbf{Persistencia del desempeño}: El puntaje en Saber 11 es el predictor más fuerte ($\beta = 0.88$, $t = 25.8$), lo que sugiere que las brechas se originan antes de la educación superior.

\item \textbf{No se encontraron diferencias por género} ($p = 0.118$).
\end{enumerate}
\end{block}
\end{frame}

\begin{frame}{Implicaciones para política educativa}
\begin{alertblock}{Recomendaciones}
\begin{enumerate}
\item \textbf{Nivelación al ingreso}: Implementar cursos de nivelación en inglés para estudiantes de estratos 1-2 y de IES públicas.

\item \textbf{Intervención temprana}: Dado que el desempeño en Saber 11 es altamente predictivo, las políticas deben enfocarse en mejorar la enseñanza del inglés en secundaria, no solo en educación superior.

\item \textbf{Equidad en recursos}: Las IES públicas requieren mayor inversión en recursos para la enseñanza del inglés (laboratorios, profesores nativos, materiales).

\item \textbf{Becas para inmersión}: Expandir programas de becas para intercambios internacionales, priorizando estudiantes de estratos bajos.

\item \textbf{Monitoreo continuo}: Evaluar periódicamente el progreso en inglés durante la carrera, no solo al final (Saber Pro).
\end{enumerate}
\end{alertblock}
\end{frame}

\begin{frame}{Limitaciones del análisis}
\begin{block}{Limitaciones metodológicas}
\begin{itemize}
\item \textbf{Causalidad vs correlación}: Los coeficientes representan asociaciones, no efectos causales. No podemos afirmar que estudiar en IES privada "causa" mayor puntaje (puede haber selección).

\item \textbf{Variables omitidas}: No controlamos por calidad del profesor, horas de estudio, motivación, experiencia previa con el inglés fuera del colegio, etc.

\item \textbf{Heterocedasticidad}: Aunque usamos errores robustos, la heterocedasticidad indica que el modelo ajusta mejor a algunos grupos que a otros.

\item \textbf{Generalización}: Los resultados son específicos para NI en Colombia en 2023. Pueden no aplicar a otros programas o países.
\end{itemize}
\end{block}

\vspace{0.3cm}
\begin{exampleblock}{Futuras investigaciones}
\begin{itemize}
\item Estudios experimentales (ej: asignación aleatoria a programas de nivelación)
\item Análisis longitudinales (seguimiento de cohortes)
\item Incorporar variables de proceso (horas de clase, metodologías)
\end{itemize}
\end{exampleblock}
\end{frame}

% ============================================================================
\section{Retroalimentación PS8}
% ============================================================================

\begin{frame}{Retroalimentación Problem Set 8}
\begin{block}{Errores comunes observados}
\begin{enumerate}
\item \textbf{No reportar diagnósticos}: Estimar OLS sin verificar supuestos.
  \begin{itemize}
  \item Solución: SIEMPRE hacer \code{plot(modelo)} y Breusch-Pagan
  \end{itemize}

\item \textbf{Interpretar coeficientes sin controles}: Decir "el estrato causa X" sin controlar por confundidores.
  \begin{itemize}
  \item Solución: usar lenguaje de asociación ("se asocia con"), no causal
  \end{itemize}

\item \textbf{Ignorar errores robustos}: Reportar errores estándar clásicos ante heterocedasticidad.
  \begin{itemize}
  \item Solución: usar \code{vcovHC} cuando Breusch-Pagan rechaza $H_0$
  \end{itemize}

\item \textbf{No reportar tamaño del efecto}: Enfocarse solo en p-valores.
  \begin{itemize}
  \item Solución: reportar coeficientes, IC, y $R^2$
  \end{itemize}

\item \textbf{Gráficos sin etiquetas}: Ejes sin nombres, sin título.
  \begin{itemize}
  \item Solución: usar \code{labs()} en ggplot
  \end{itemize}
\end{enumerate}
\end{block}
\end{frame}

\begin{frame}{Buenas prácticas observadas}
\begin{exampleblock}{Ejemplos de excelencia}
\begin{itemize}
\item \textbf{Análisis exploratorio robusto}: Desagregar por múltiples grupos, graficar antes de modelar.

\item \textbf{Comparación de modelos}: Reportar modelo simple y múltiple, discutir OVB.

\item \textbf{Análisis de sensibilidad}: Re-estimar sin observaciones influyentes, reportar si cambian conclusiones.

\item \textbf{Visualizaciones efectivas}: Forest plots, boxplots comparativos, scatterplots con línea de regresión.

\item \textbf{Interpretación contextualizada}: Conectar resultados estadísticos con implicaciones para NI (ej: "Esta brecha equivale a medio nivel MCER").

\item \textbf{Código reproducible}: Scripts bien comentados, con estructura clara (carga $\to$ limpieza $\to$ EDA $\to$ modelos).
\end{itemize}
\end{exampleblock}

\vspace{0.3cm}
\small
\textcolor{javeriana}{Recordatorio}: El código es parte del informe. Debe ser legible y reproducible.
\end{frame}

\begin{frame}{Rúbrica para el proyecto final}
\small
\begin{center}
\renewcommand{\arraystretch}{1.3}
\begin{tabular}{lp{6cm}r}
\toprule
\textbf{Criterio} & \textbf{Descripción} & \textbf{Puntos} \\
\midrule
Pregunta de investigación & Clara, medible, relevante & 10 \\
Limpieza de datos & Tratamiento de NA, valores atípicos, creación de variables & 10 \\
EDA & Estadísticos descriptivos, visualizaciones, desagregación & 15 \\
Inferencia & IC, pruebas de hipótesis, interpretación & 15 \\
Modelo de regresión & Especificación, interpretación, comparación & 20 \\
Diagnósticos & Gráficos, tests, correcciones (errores robustos) & 15 \\
Conclusiones & Síntesis, implicaciones, limitaciones & 10 \\
Código reproducible & Scripts claros, comentados, reproducibles & 5 \\
\midrule
\textbf{Total} &  & \textbf{100} \\
\bottomrule
\end{tabular}
\end{center}

\vspace{0.3cm}
\textbf{Entrega}: Informe en PDF ($\leq$15 páginas) + código R (.R o .Rmd) + datos.\\
\textbf{Fecha límite}: Viernes 20 de noviembre, 23:59h (Moodle).
\end{frame}

% ============================================================================
\section{Resumen}
% ============================================================================

\begin{frame}{El pipeline estadístico: repaso}
\begin{center}
\begin{tikzpicture}[
  node distance=0.6cm,
  box/.style={rectangle, draw=javeriana, fill=javeriana!10, thick, minimum width=3cm, minimum height=0.6cm, rounded corners, font=\small}
]
\node[box] (p1) {1. Pregunta clara};
\node[box, below=of p1] (p2) {2. Datos limpios};
\node[box, below=of p2] (p3) {3. EDA exhaustivo};
\node[box, below=of p3] (p4) {4. IC para incertidumbre};
\node[box, below=of p4] (p5) {5. Pruebas para decisiones};
\node[box, below=of p5] (p6) {6. Modelo para relaciones};
\node[box, below=of p6] (p7) {7. Diagnósticos para validez};
\node[box, below=of p7] (p8) {8. Conclusiones contextualizadas};

\draw[->, thick, javeriana] (p1) -- (p2);
\draw[->, thick, javeriana] (p2) -- (p3);
\draw[->, thick, javeriana] (p3) -- (p4);
\draw[->, thick, javeriana] (p4) -- (p5);
\draw[->, thick, javeriana] (p5) -- (p6);
\draw[->, thick, javeriana] (p6) -- (p7);
\draw[->, thick, javeriana] (p7) -- (p8);
\end{tikzpicture}
\end{center}
\end{frame}

\begin{frame}{Lecciones clave}
\begin{alertblock}{Principios del análisis estadístico riguroso}
\begin{enumerate}
\item \textbf{Explorar antes de modelar}: Los datos deben "hablarte" antes de imponer un modelo.

\item \textbf{Controlar confundidores}: Un modelo simple puede ser engañoso (OVB).

\item \textbf{Verificar supuestos}: OLS sin diagnósticos es como conducir sin frenos.

\item \textbf{Reportar incertidumbre}: Un coeficiente sin IC o error estándar es incompleto.

\item \textbf{Interpretar en contexto}: Los números cobran sentido al conectarlos con la pregunta original.

\item \textbf{Reconocer limitaciones}: Un buen análisis admite lo que NO puede concluir.

\item \textbf{Ser reproducible}: Tu código es tu receta. Otros (y tú en el futuro) deben poder replicar tus resultados.
\end{enumerate}
\end{alertblock}
\end{frame}

\begin{frame}{Próxima sesión: repaso final}
\begin{block}{Sesión 13: Repaso Final y Preparación para el Examen}
\begin{itemize}
\item Mapa conceptual del curso (4 unidades)
\item Repaso de fórmulas clave (IC, pruebas, regresión)
\item Problemas tipo examen
\item Estrategias para el examen final
\item Instrucciones del proyecto final
\end{itemize}
\end{block}

\vspace{0.3cm}
\begin{alertblock}{Para la próxima clase}
\textbf{Traer:}
\begin{itemize}
\item Todas tus dudas del semestre
\item Hoja de fórmulas (draft para el examen)
\item Calculadora
\end{itemize}
\end{alertblock}
\end{frame}

\begin{frame}{}
\begin{center}
\Huge ¡Gracias!

\vspace{1cm}
\Large Preguntas y discusión

\vspace{1.5cm}
\normalsize
\textbf{Contacto:}\\
jaime.polanco@javeriana.edu.co\\
Horario de atención: Miércoles 14:00-16:00
\end{center}
\end{frame}

\end{document}
